\family{Meliaceae}{Mahogany family}
\label{fam:meliaceae}

Trees or shrubs, the bark and wood often scented. Leaves alternate,
usually pinnate, without stipules. Flowers small, 4- or 5-merous, the
stamens joined into a tube. Fruit a capsule, a berry or a drupe; in
the timber genera treated here a woody capsule with winged seeds.
Generic limits follow \citet{pennington1975}. Three species in the
Arami Uplands, all valuable timber trees.

\section*{Key to the species}
\begin{florakey}
	\lead{1a}{Capsule cylindrical, pendulous, splitting from the base into
		5 valves; slash with a cedar-like scent}{\taxon{Entandrophragma
			cylindricum}}
	\lead{1b}{Capsule globose, erect, opening from the apex into 4 or 5
		valves}{2}
	\lead{2a}{Leaflets broadly obovate, up to 20~cm long; capsule
		6--8~cm across; forest tree}{\taxon{Khaya grandifoliola}}
	\lead{2b}{Leaflets oblong-elliptic, up to 10~cm long; capsule
		4--6~cm across; savanna tree}{\taxon{Khaya senegalensis}}
\end{florakey}

\species{Entandrophragma cylindricum}{(Sprague) Sprague}
\label{sp:MEL01}

\field{Synonyms} \sciname{Pseudocedrela cylindrica}{Sprague}

\field{Local and trade names} Trade: Sapele.

\field{Status} Native.

\field{Description} Deciduous tree to 50~m with short, blunt
buttresses. Bark grey-brown, flaking in small round scales; slash with
a cedar-like scent. Leaves paripinnate with 5--9 pairs of leaflets.
Flowers small, greenish-yellow. Capsule cylindrical, pendulous,
6--10~cm long, splitting from the base into 5 valves; seeds winged.

\field{Phenology} Flowering March to May; fruiting November to
January.

\field{Habitat} Lowland semi-deciduous forest.

\field{Uses} A reddish-brown timber with a striped figure, used for
veneer, furniture and joinery. Heavy logging has made large trees
scarce \citep{hawthorne2006,louppe2008}.

\field{Specimens examined} Ireke Ferry (D4), 360~m, fr., 10~Jan 1934, \specimen{Harding}{1059}{OKF};
Ireke Ferry (D4), 360~m, st., 20~Sep 1934, \specimen{Harding}{1209}{OKF};
Ireke Ferry (D4), 360~m, fr., 14~Dec 1936, \specimen{Harding}{1598}{OKF};
Okeri (C4), 420~m, fr., 17~Dec 1938, \specimen{Harding}{2055}{OKF};
Ireke Ferry (D4), 360~m, fr., 19~Jan 2019, \specimen{Alade}{145}{AUH};
Olowo Stream (D3), 440~m, fr., 2~Jan 2024, \specimen{Alade}{840}{AUH};
Olowo Stream (D3), 440~m, fr., 3~Dec 2024, \specimen{Alade}{938}{AUH}.

\field{Distribution} Grid squares C4, D3, D4.

\begin{center}
\distmap{C/4, D/3, D/4}
\end{center}

\species{Khaya grandifoliola}{C.DC.}
\label{sp:MEL02}

\field{Local and trade names} Yoruba: Oganwo. Trade: African
mahogany.

\field{Status} Native.

\field{Description} Deciduous tree to 40~m with buttresses. Bark
grey-brown, flaking in scales; slash red. Leaves paripinnate with
2--5 pairs of broadly obovate leaflets up to 20~cm long. Flowers
small, cream. Capsule erect, globose, 6--8~cm across, opening into 4
or 5 valves; seeds winged.

\field{Phenology} Flowering March to May; fruiting January to March.

\field{Habitat} Lowland semi-deciduous forest and forest--savanna
mosaic.

\field{Uses} A valuable mahogany timber for furniture and joinery. The
bitter bark is used in local medicine \citep{keay1989,louppe2008}.

\field{Specimens examined} Ireke Ferry (D4), 360~m, fr., 28~Jan 1936, \specimen{Harding}{1391}{OKF};
Olowo Stream (D3), 440~m, st., 6~Nov 1937, \specimen{Harding}{1805}{OKF};
Ireke Ferry (D4), 360~m, fr., 20~Mar 1938, \specimen{Harding}{1868}{OKF};
Ewe Stream gallery forest (B3), 580~m, fr., 10~Feb 1961, \specimen{Olaniyan}{293}{SVN};
Ewe Stream gallery forest (B3), 580~m, fr., 8~Feb 1966, \specimen{Olaniyan}{774}{SVN};
Okeri (C4), 420~m, fr., 15~Mar 2019, \specimen{Alade}{189}{AUH};
Agbede (C2), 540~m, fr., 26~Mar 2020, \specimen{Alade}{263}{AUH}.

\field{Distribution} Grid squares B3, C2, C4, D3, D4.

\begin{center}
\distmap{B/3, C/2, C/4, D/3, D/4}
\end{center}

\species{Khaya senegalensis}{(Desr.) A.Juss.}
\label{sp:MEL03}

\field{Synonyms} \sciname{Swietenia senegalensis}{Desr.}

\field{Local and trade names} Hausa: Madachi. Trade: Dry-zone
mahogany.

\field{Status} Native.

\field{Description} Tree to 30~m with a dense crown. Bark dark grey,
scaly. Leaves paripinnate with 3--6 pairs of oblong-elliptic leaflets
up to 10~cm long. Flowers small, cream. Capsule erect, globose,
4--6~cm across, opening into 4 valves; seeds winged.

\field{Phenology} Flowering March to May; fruiting January to
February.

\field{Habitat} Guinea savanna woodland, especially along
watercourses; also planted as a shade and avenue tree.

\field{Uses} Timber for furniture and building. The very bitter bark
is used in local medicine, and the branches are lopped for fodder in
the dry season \citep{arbonnier2004,orwa2009}.

\field{Specimens examined} Aramiwo Escarpment, summit (B1), 900~m, fr., 22~Feb 1962, \specimen{Olaniyan}{360}{SVN};
Idofin (A3), 760~m, fr., 14~Jan 1963, \specimen{Olaniyan}{515}{SVN};
Aramiwo Escarpment, summit (B1), 900~m, fr., 26~Jan 1967, \specimen{Olaniyan}{878}{SVN};
Aramiwo Escarpment, summit (B1), 900~m, fr., 13~Feb 1967, \specimen{Olaniyan}{909}{SVN};
Idofin (A3), 760~m, fl., 10~Apr 2019, \specimen{Alade}{207}{AUH};
Aramiwo Escarpment, summit (B1), 900~m, fl., 6~Mar 2023, \specimen{Alade}{685}{AUH};
Asiri Crossing (A2), 740~m, fr., 22~Jan 2025, \specimen{Alade}{974}{AUH}.

\field{Distribution} Grid squares A2, A3, B1.

\begin{center}
\distmap{A/2, A/3, B/1}
\end{center}
